% ============================================================
%  preface_ai.tex  —  Preface: Use of AI
%  Referenced by olowiane_dzieci_analiza.tex via % ============================================================
%  preface_ai.tex  —  Preface: Use of AI
%  Referenced by olowiane_dzieci_analiza.tex via % ============================================================
%  preface_ai.tex  —  Preface: Use of AI
%  Referenced by olowiane_dzieci_analiza.tex via % ============================================================
%  preface_ai.tex  —  Preface: Use of AI
%  Referenced by olowiane_dzieci_analiza.tex via \input{preface_ai}
% ============================================================

\chapter*{A Note on Method: Use of Artificial Intelligence}
\addcontentsline{toc}{chapter}{A Note on Method: Use of AI}
\markboth{A Note on Method}{}

This document was produced with the assistance of a large language
model (LLM) and is shared openly in the hope that others may adapt
the method for their own language learning.

\section*{The Learning Context}

The author is a native English speaker with a Polish wife, Polish
in-laws, and a recently acquired apartment in Kraków. After more than
twenty years of incidental exposure to Polish, conversational ability
had plateaued at a level sufficient for simple transactions but
inadequate for real two-way communication. The goal of this project
is to break through that plateau by engaging deeply with a Polish-language
book that is genuinely interesting, not available in English translation,
and whose subject matter --- a lead poisoning epidemic in industrial
Silesia --- connects directly to the author's professional background
in biomedical engineering.

The author has ADHD and dyslexia/dysorthographia, and shares this
openly in the hope of reducing the stigma that still surrounds these
conditions --- particularly in academic and language-learning contexts,
where they are often misread as indicators of low ability or
insufficient effort. \LaTeX{} is used throughout as an accessibility
aid: its high-quality typesetting, structured navigation,
wheat-coloured page background, and clean separation between source
editing and rendered output all materially reduce visual fatigue
compared to word processors. This is noted here because the same
approach may benefit other learners with similar profiles.

Neuropsychological assessment has identified two features of the
author's cognitive profile that are directly relevant to this
project and worth noting for other learners who may recognise
themselves in them:

\begin{description}

  \item[Reduced working memory] Associated with ADHD, this makes
    rote memorisation of vocabulary lists or paradigm tables
    effortful and largely ineffective as a learning strategy.
    Accordingly, no attempt is made to memorise the material in
    this document. The goal is repeated, pleasurable exposure ---
    the same words and constructions reappearing in new contexts
    until they become familiar without deliberate effort.
    The analyses are written to be enjoyable to read in their own
    right, not to be studied.

  \item[Strong pattern recognition and rule derivation] Identified
    as a compensating strength, this is the capacity to infer
    underlying rules and structures from examples, and to
    generalise them rapidly once the pattern is seen. This is
    precisely what grammatical analysis rewards: rather than
    memorising that \textit{na cmentarzu} means ``at the cemetery,''
    the learner recognises the locative pattern, derives the rule,
    and thereafter reads all locative constructions with growing
    confidence --- without having memorised a single paradigm table.
    The Latin background established this kind of structural
    awareness early; Polish is being acquired in the same way.

\end{description}

These two features together suggest a learning profile that is
poorly served by conventional language instruction (which relies
heavily on working memory and rote practice) but well served by
immersive, contextual, analysis-rich reading of genuinely engaging
material. The method described in this document was developed
organically around exactly this profile.

\section*{The AI Tool}

All phrase-by-phrase translation, grammatical analysis, and literary
commentary in this document was generated interactively using
\textbf{Claude}, an AI assistant developed by
\href{https://www.anthropic.com}{Anthropic}. Specifically, the model
used was \textbf{Claude Sonnet 4.6}
(model string: \texttt{claude-sonnet-4-6}), accessed via the
\href{https://claude.ai}{Claude.ai} web interface.

The \LaTeX{} document itself --- its structure, preamble, custom
commands, grammar reference, and chapter files --- was also drafted
and maintained by the model in response to the author's direction,
updated periodically as new material accumulated.

\section*{The Method}

The workflow is simple and requires no technical expertise beyond
basic \LaTeX{} compilation:

\begin{enumerate}

  \item The learner opens the source book in an e-reader application.
  Since the book is protected by digital rights management (DRM),
  text cannot be copied directly. Instead, the learner takes a
  screenshot of one sentence or short passage at a time.

  \item The screenshot is pasted into a conversation with Claude.
  The model reads the Polish text from the image and provides:
  \begin{itemize}
    \item a full English translation of the passage;
    \item a word-by-word or phrase-by-phrase breakdown with
      vocabulary glosses;
    \item grammatical analysis --- case endings, verb aspect,
      participial constructions, and other relevant features;
    \item literary and contextual commentary where the author's
      choices are particularly interesting.
  \end{itemize}

  \item The learner reads the analysis, asks follow-up questions
  (about vocabulary, grammar, or usage), and continues to the
  next passage at their own pace. There is no pressure to
  memorise --- the goal is absorption through engaged reading.

  \item Periodically, at natural section breaks, the accumulated
  analyses are consolidated into the \LaTeX{} document. The model
  drafts the \LaTeX{} source; the learner compiles it locally
  and reviews the output.

  \item When the AI session approaches its context window limit,
  a new session is started. The main \texttt{.tex} file --- which
  contains a detailed comment block describing the project,
  learner profile, and working method --- is uploaded to the new
  session as a handoff document, allowing work to resume
  seamlessly.

\end{enumerate}

\section*{Depth of Analysis}

The depth of analysis is intentionally variable. Some sentences
are unpacked in full detail because the grammar or vocabulary is
instructive. Others are translated more briefly because the content
is transitional or the constructions are already familiar. The
learner directs the pace: tangential questions about grammar or
usage are encouraged at any point, and the conversation moves
between close analysis and broader literary observation as the
text warrants.

This variability is a feature, not a deficiency. It mirrors the
way a knowledgeable reading companion would engage with a text ---
pausing where things are interesting, moving on where they are not.

\section*{Suitability for Other Languages and Learners}

The method is not specific to Polish. It is applicable to any
language for which a sufficiently capable LLM exists, and to any
learner who has a genuine motivation to read a specific text in
that language. The key conditions that make it work are:

\begin{itemize}
  \item a source text that the learner \textit{actually wants to read}
    --- ideally one unavailable in their native language, removing
    the option of simply reading a translation;
  \item willingness to proceed slowly and to treat the analysis as
    enjoyable in its own right, not merely as a means to an end;
  \item basic familiarity with \LaTeX{} if the learner wishes to
    maintain a compiled document, though the conversational
    analysis alone is valuable without any document production.
\end{itemize}

The method works particularly well for learners with existing
metalinguistic awareness --- those who have studied another
inflected language (Latin, Russian, German, classical Greek) and
already understand concepts like case, aspect, and agreement
abstractly. For such learners, the AI analysis illuminates the
\textit{Polish instantiation} of familiar concepts rather than
introducing entirely new ideas from scratch.

\section*{Transparency and Limitations}

The analyses in this document reflect the capabilities of the model
at the time of writing and should be treated as a knowledgeable
first-pass rather than authoritative scholarship. Occasional errors
in grammatical analysis or nuance are possible. Where the author
has Polish-speaking family or other expert resources available,
those remain the final arbiter of correctness.

The quoted Polish passages in this document are brief and used
solely for educational commentary. They do not constitute
reproduction of the source text and are consistent with fair use
principles for personal, non-commercial educational purposes.
The document is shared openly but the source book remains under
its original copyright.

\vspace{1.5em}
\noindent The full project, including all \LaTeX{} source files,
is available at:\\
\href{https://github.com/}{[GitHub repository URL to be added]}

\vspace{0.5em}
\noindent Readers are welcome to adapt the structure, commands,
and approach for their own language learning projects.

% ============================================================

\chapter*{A Note on Method: Use of Artificial Intelligence}
\addcontentsline{toc}{chapter}{A Note on Method: Use of AI}
\markboth{A Note on Method}{}

This document was produced with the assistance of a large language
model (LLM) and is shared openly in the hope that others may adapt
the method for their own language learning.

\section*{The Learning Context}

The author is a native English speaker with a Polish wife, Polish
in-laws, and a recently acquired apartment in Kraków. After more than
twenty years of incidental exposure to Polish, conversational ability
had plateaued at a level sufficient for simple transactions but
inadequate for real two-way communication. The goal of this project
is to break through that plateau by engaging deeply with a Polish-language
book that is genuinely interesting, not available in English translation,
and whose subject matter --- a lead poisoning epidemic in industrial
Silesia --- connects directly to the author's professional background
in biomedical engineering.

The author has ADHD and dyslexia/dysorthographia, and shares this
openly in the hope of reducing the stigma that still surrounds these
conditions --- particularly in academic and language-learning contexts,
where they are often misread as indicators of low ability or
insufficient effort. \LaTeX{} is used throughout as an accessibility
aid: its high-quality typesetting, structured navigation,
wheat-coloured page background, and clean separation between source
editing and rendered output all materially reduce visual fatigue
compared to word processors. This is noted here because the same
approach may benefit other learners with similar profiles.

Neuropsychological assessment has identified two features of the
author's cognitive profile that are directly relevant to this
project and worth noting for other learners who may recognise
themselves in them:

\begin{description}

  \item[Reduced working memory] Associated with ADHD, this makes
    rote memorisation of vocabulary lists or paradigm tables
    effortful and largely ineffective as a learning strategy.
    Accordingly, no attempt is made to memorise the material in
    this document. The goal is repeated, pleasurable exposure ---
    the same words and constructions reappearing in new contexts
    until they become familiar without deliberate effort.
    The analyses are written to be enjoyable to read in their own
    right, not to be studied.

  \item[Strong pattern recognition and rule derivation] Identified
    as a compensating strength, this is the capacity to infer
    underlying rules and structures from examples, and to
    generalise them rapidly once the pattern is seen. This is
    precisely what grammatical analysis rewards: rather than
    memorising that \textit{na cmentarzu} means ``at the cemetery,''
    the learner recognises the locative pattern, derives the rule,
    and thereafter reads all locative constructions with growing
    confidence --- without having memorised a single paradigm table.
    The Latin background established this kind of structural
    awareness early; Polish is being acquired in the same way.

\end{description}

These two features together suggest a learning profile that is
poorly served by conventional language instruction (which relies
heavily on working memory and rote practice) but well served by
immersive, contextual, analysis-rich reading of genuinely engaging
material. The method described in this document was developed
organically around exactly this profile.

\section*{The AI Tool}

All phrase-by-phrase translation, grammatical analysis, and literary
commentary in this document was generated interactively using
\textbf{Claude}, an AI assistant developed by
\href{https://www.anthropic.com}{Anthropic}. Specifically, the model
used was \textbf{Claude Sonnet 4.6}
(model string: \texttt{claude-sonnet-4-6}), accessed via the
\href{https://claude.ai}{Claude.ai} web interface.

The \LaTeX{} document itself --- its structure, preamble, custom
commands, grammar reference, and chapter files --- was also drafted
and maintained by the model in response to the author's direction,
updated periodically as new material accumulated.

\section*{The Method}

The workflow is simple and requires no technical expertise beyond
basic \LaTeX{} compilation:

\begin{enumerate}

  \item The learner opens the source book in an e-reader application.
  Since the book is protected by digital rights management (DRM),
  text cannot be copied directly. Instead, the learner takes a
  screenshot of one sentence or short passage at a time.

  \item The screenshot is pasted into a conversation with Claude.
  The model reads the Polish text from the image and provides:
  \begin{itemize}
    \item a full English translation of the passage;
    \item a word-by-word or phrase-by-phrase breakdown with
      vocabulary glosses;
    \item grammatical analysis --- case endings, verb aspect,
      participial constructions, and other relevant features;
    \item literary and contextual commentary where the author's
      choices are particularly interesting.
  \end{itemize}

  \item The learner reads the analysis, asks follow-up questions
  (about vocabulary, grammar, or usage), and continues to the
  next passage at their own pace. There is no pressure to
  memorise --- the goal is absorption through engaged reading.

  \item Periodically, at natural section breaks, the accumulated
  analyses are consolidated into the \LaTeX{} document. The model
  drafts the \LaTeX{} source; the learner compiles it locally
  and reviews the output.

  \item When the AI session approaches its context window limit,
  a new session is started. The main \texttt{.tex} file --- which
  contains a detailed comment block describing the project,
  learner profile, and working method --- is uploaded to the new
  session as a handoff document, allowing work to resume
  seamlessly.

\end{enumerate}

\section*{Depth of Analysis}

The depth of analysis is intentionally variable. Some sentences
are unpacked in full detail because the grammar or vocabulary is
instructive. Others are translated more briefly because the content
is transitional or the constructions are already familiar. The
learner directs the pace: tangential questions about grammar or
usage are encouraged at any point, and the conversation moves
between close analysis and broader literary observation as the
text warrants.

This variability is a feature, not a deficiency. It mirrors the
way a knowledgeable reading companion would engage with a text ---
pausing where things are interesting, moving on where they are not.

\section*{Suitability for Other Languages and Learners}

The method is not specific to Polish. It is applicable to any
language for which a sufficiently capable LLM exists, and to any
learner who has a genuine motivation to read a specific text in
that language. The key conditions that make it work are:

\begin{itemize}
  \item a source text that the learner \textit{actually wants to read}
    --- ideally one unavailable in their native language, removing
    the option of simply reading a translation;
  \item willingness to proceed slowly and to treat the analysis as
    enjoyable in its own right, not merely as a means to an end;
  \item basic familiarity with \LaTeX{} if the learner wishes to
    maintain a compiled document, though the conversational
    analysis alone is valuable without any document production.
\end{itemize}

The method works particularly well for learners with existing
metalinguistic awareness --- those who have studied another
inflected language (Latin, Russian, German, classical Greek) and
already understand concepts like case, aspect, and agreement
abstractly. For such learners, the AI analysis illuminates the
\textit{Polish instantiation} of familiar concepts rather than
introducing entirely new ideas from scratch.

\section*{Transparency and Limitations}

The analyses in this document reflect the capabilities of the model
at the time of writing and should be treated as a knowledgeable
first-pass rather than authoritative scholarship. Occasional errors
in grammatical analysis or nuance are possible. Where the author
has Polish-speaking family or other expert resources available,
those remain the final arbiter of correctness.

The quoted Polish passages in this document are brief and used
solely for educational commentary. They do not constitute
reproduction of the source text and are consistent with fair use
principles for personal, non-commercial educational purposes.
The document is shared openly but the source book remains under
its original copyright.

\vspace{1.5em}
\noindent The full project, including all \LaTeX{} source files,
is available at:\\
\href{https://github.com/}{[GitHub repository URL to be added]}

\vspace{0.5em}
\noindent Readers are welcome to adapt the structure, commands,
and approach for their own language learning projects.

% ============================================================

\chapter*{A Note on Method: Use of Artificial Intelligence}
\addcontentsline{toc}{chapter}{A Note on Method: Use of AI}
\markboth{A Note on Method}{}

This document was produced with the assistance of a large language
model (LLM) and is shared openly in the hope that others may adapt
the method for their own language learning.

\section*{The Learning Context}

The author is a native English speaker with a Polish wife, Polish
in-laws, and a recently acquired apartment in Kraków. After more than
twenty years of incidental exposure to Polish, conversational ability
had plateaued at a level sufficient for simple transactions but
inadequate for real two-way communication. The goal of this project
is to break through that plateau by engaging deeply with a Polish-language
book that is genuinely interesting, not available in English translation,
and whose subject matter --- a lead poisoning epidemic in industrial
Silesia --- connects directly to the author's professional background
in biomedical engineering.

The author has ADHD and dyslexia/dysorthographia, and shares this
openly in the hope of reducing the stigma that still surrounds these
conditions --- particularly in academic and language-learning contexts,
where they are often misread as indicators of low ability or
insufficient effort. \LaTeX{} is used throughout as an accessibility
aid: its high-quality typesetting, structured navigation,
wheat-coloured page background, and clean separation between source
editing and rendered output all materially reduce visual fatigue
compared to word processors. This is noted here because the same
approach may benefit other learners with similar profiles.

Neuropsychological assessment has identified two features of the
author's cognitive profile that are directly relevant to this
project and worth noting for other learners who may recognise
themselves in them:

\begin{description}

  \item[Reduced working memory] Associated with ADHD, this makes
    rote memorisation of vocabulary lists or paradigm tables
    effortful and largely ineffective as a learning strategy.
    Accordingly, no attempt is made to memorise the material in
    this document. The goal is repeated, pleasurable exposure ---
    the same words and constructions reappearing in new contexts
    until they become familiar without deliberate effort.
    The analyses are written to be enjoyable to read in their own
    right, not to be studied.

  \item[Strong pattern recognition and rule derivation] Identified
    as a compensating strength, this is the capacity to infer
    underlying rules and structures from examples, and to
    generalise them rapidly once the pattern is seen. This is
    precisely what grammatical analysis rewards: rather than
    memorising that \textit{na cmentarzu} means ``at the cemetery,''
    the learner recognises the locative pattern, derives the rule,
    and thereafter reads all locative constructions with growing
    confidence --- without having memorised a single paradigm table.
    The Latin background established this kind of structural
    awareness early; Polish is being acquired in the same way.

\end{description}

These two features together suggest a learning profile that is
poorly served by conventional language instruction (which relies
heavily on working memory and rote practice) but well served by
immersive, contextual, analysis-rich reading of genuinely engaging
material. The method described in this document was developed
organically around exactly this profile.

\section*{The AI Tool}

All phrase-by-phrase translation, grammatical analysis, and literary
commentary in this document was generated interactively using
\textbf{Claude}, an AI assistant developed by
\href{https://www.anthropic.com}{Anthropic}. Specifically, the model
used was \textbf{Claude Sonnet 4.6}
(model string: \texttt{claude-sonnet-4-6}), accessed via the
\href{https://claude.ai}{Claude.ai} web interface.

The \LaTeX{} document itself --- its structure, preamble, custom
commands, grammar reference, and chapter files --- was also drafted
and maintained by the model in response to the author's direction,
updated periodically as new material accumulated.

\section*{The Method}

The workflow is simple and requires no technical expertise beyond
basic \LaTeX{} compilation:

\begin{enumerate}

  \item The learner opens the source book in an e-reader application.
  Since the book is protected by digital rights management (DRM),
  text cannot be copied directly. Instead, the learner takes a
  screenshot of one sentence or short passage at a time.

  \item The screenshot is pasted into a conversation with Claude.
  The model reads the Polish text from the image and provides:
  \begin{itemize}
    \item a full English translation of the passage;
    \item a word-by-word or phrase-by-phrase breakdown with
      vocabulary glosses;
    \item grammatical analysis --- case endings, verb aspect,
      participial constructions, and other relevant features;
    \item literary and contextual commentary where the author's
      choices are particularly interesting.
  \end{itemize}

  \item The learner reads the analysis, asks follow-up questions
  (about vocabulary, grammar, or usage), and continues to the
  next passage at their own pace. There is no pressure to
  memorise --- the goal is absorption through engaged reading.

  \item Periodically, at natural section breaks, the accumulated
  analyses are consolidated into the \LaTeX{} document. The model
  drafts the \LaTeX{} source; the learner compiles it locally
  and reviews the output.

  \item When the AI session approaches its context window limit,
  a new session is started. The main \texttt{.tex} file --- which
  contains a detailed comment block describing the project,
  learner profile, and working method --- is uploaded to the new
  session as a handoff document, allowing work to resume
  seamlessly.

\end{enumerate}

\section*{Depth of Analysis}

The depth of analysis is intentionally variable. Some sentences
are unpacked in full detail because the grammar or vocabulary is
instructive. Others are translated more briefly because the content
is transitional or the constructions are already familiar. The
learner directs the pace: tangential questions about grammar or
usage are encouraged at any point, and the conversation moves
between close analysis and broader literary observation as the
text warrants.

This variability is a feature, not a deficiency. It mirrors the
way a knowledgeable reading companion would engage with a text ---
pausing where things are interesting, moving on where they are not.

\section*{Suitability for Other Languages and Learners}

The method is not specific to Polish. It is applicable to any
language for which a sufficiently capable LLM exists, and to any
learner who has a genuine motivation to read a specific text in
that language. The key conditions that make it work are:

\begin{itemize}
  \item a source text that the learner \textit{actually wants to read}
    --- ideally one unavailable in their native language, removing
    the option of simply reading a translation;
  \item willingness to proceed slowly and to treat the analysis as
    enjoyable in its own right, not merely as a means to an end;
  \item basic familiarity with \LaTeX{} if the learner wishes to
    maintain a compiled document, though the conversational
    analysis alone is valuable without any document production.
\end{itemize}

The method works particularly well for learners with existing
metalinguistic awareness --- those who have studied another
inflected language (Latin, Russian, German, classical Greek) and
already understand concepts like case, aspect, and agreement
abstractly. For such learners, the AI analysis illuminates the
\textit{Polish instantiation} of familiar concepts rather than
introducing entirely new ideas from scratch.

\section*{Transparency and Limitations}

The analyses in this document reflect the capabilities of the model
at the time of writing and should be treated as a knowledgeable
first-pass rather than authoritative scholarship. Occasional errors
in grammatical analysis or nuance are possible. Where the author
has Polish-speaking family or other expert resources available,
those remain the final arbiter of correctness.

The quoted Polish passages in this document are brief and used
solely for educational commentary. They do not constitute
reproduction of the source text and are consistent with fair use
principles for personal, non-commercial educational purposes.
The document is shared openly but the source book remains under
its original copyright.

\vspace{1.5em}
\noindent The full project, including all \LaTeX{} source files,
is available at:\\
\href{https://github.com/}{[GitHub repository URL to be added]}

\vspace{0.5em}
\noindent Readers are welcome to adapt the structure, commands,
and approach for their own language learning projects.

% ============================================================

\chapter*{A Note on Method: Use of Artificial Intelligence}
\addcontentsline{toc}{chapter}{A Note on Method: Use of AI}
\markboth{A Note on Method}{}

This document was produced with the assistance of a large language
model (LLM) and is shared openly in the hope that others may adapt
the method for their own language learning.

\section*{The Learning Context}

The author is a native English speaker with a Polish wife, Polish
in-laws, and a recently acquired apartment in Kraków. After more than
twenty years of incidental exposure to Polish, conversational ability
had plateaued at a level sufficient for simple transactions but
inadequate for real two-way communication. The goal of this project
is to break through that plateau by engaging deeply with a Polish-language
book that is genuinely interesting, not available in English translation,
and whose subject matter --- a lead poisoning epidemic in industrial
Silesia --- connects directly to the author's professional background
in biomedical engineering.

The author has ADHD and dyslexia/dysorthographia, and shares this
openly in the hope of reducing the stigma that still surrounds these
conditions --- particularly in academic and language-learning contexts,
where they are often misread as indicators of low ability or
insufficient effort. \LaTeX{} is used throughout as an accessibility
aid: its high-quality typesetting, structured navigation,
wheat-coloured page background, and clean separation between source
editing and rendered output all materially reduce visual fatigue
compared to word processors. This is noted here because the same
approach may benefit other learners with similar profiles.

Neuropsychological assessment has identified two features of the
author's cognitive profile that are directly relevant to this
project and worth noting for other learners who may recognise
themselves in them:

\begin{description}

  \item[Reduced working memory] Associated with ADHD, this makes
    rote memorisation of vocabulary lists or paradigm tables
    effortful and largely ineffective as a learning strategy.
    Accordingly, no attempt is made to memorise the material in
    this document. The goal is repeated, pleasurable exposure ---
    the same words and constructions reappearing in new contexts
    until they become familiar without deliberate effort.
    The analyses are written to be enjoyable to read in their own
    right, not to be studied.

  \item[Strong pattern recognition and rule derivation] Identified
    as a compensating strength, this is the capacity to infer
    underlying rules and structures from examples, and to
    generalise them rapidly once the pattern is seen. This is
    precisely what grammatical analysis rewards: rather than
    memorising that \textit{na cmentarzu} means ``at the cemetery,''
    the learner recognises the locative pattern, derives the rule,
    and thereafter reads all locative constructions with growing
    confidence --- without having memorised a single paradigm table.
    The Latin background established this kind of structural
    awareness early; Polish is being acquired in the same way.

\end{description}

These two features together suggest a learning profile that is
poorly served by conventional language instruction (which relies
heavily on working memory and rote practice) but well served by
immersive, contextual, analysis-rich reading of genuinely engaging
material. The method described in this document was developed
organically around exactly this profile.

\section*{The AI Tool}

All phrase-by-phrase translation, grammatical analysis, and literary
commentary in this document was generated interactively using
\textbf{Claude}, an AI assistant developed by
\href{https://www.anthropic.com}{Anthropic}. Specifically, the model
used was \textbf{Claude Sonnet 4.6}
(model string: \texttt{claude-sonnet-4-6}), accessed via the
\href{https://claude.ai}{Claude.ai} web interface.

The \LaTeX{} document itself --- its structure, preamble, custom
commands, grammar reference, and chapter files --- was also drafted
and maintained by the model in response to the author's direction,
updated periodically as new material accumulated.

\section*{The Method}

The workflow is simple and requires no technical expertise beyond
basic \LaTeX{} compilation:

\begin{enumerate}

  \item The learner opens the source book in an e-reader application.
  Since the book is protected by digital rights management (DRM),
  text cannot be copied directly. Instead, the learner takes a
  screenshot of one sentence or short passage at a time.

  \item The screenshot is pasted into a conversation with Claude.
  The model reads the Polish text from the image and provides:
  \begin{itemize}
    \item a full English translation of the passage;
    \item a word-by-word or phrase-by-phrase breakdown with
      vocabulary glosses;
    \item grammatical analysis --- case endings, verb aspect,
      participial constructions, and other relevant features;
    \item literary and contextual commentary where the author's
      choices are particularly interesting.
  \end{itemize}

  \item The learner reads the analysis, asks follow-up questions
  (about vocabulary, grammar, or usage), and continues to the
  next passage at their own pace. There is no pressure to
  memorise --- the goal is absorption through engaged reading.

  \item Periodically, at natural section breaks, the accumulated
  analyses are consolidated into the \LaTeX{} document. The model
  drafts the \LaTeX{} source; the learner compiles it locally
  and reviews the output.

  \item When the AI session approaches its context window limit,
  a new session is started. The main \texttt{.tex} file --- which
  contains a detailed comment block describing the project,
  learner profile, and working method --- is uploaded to the new
  session as a handoff document, allowing work to resume
  seamlessly.

\end{enumerate}

\section*{Depth of Analysis}

The depth of analysis is intentionally variable. Some sentences
are unpacked in full detail because the grammar or vocabulary is
instructive. Others are translated more briefly because the content
is transitional or the constructions are already familiar. The
learner directs the pace: tangential questions about grammar or
usage are encouraged at any point, and the conversation moves
between close analysis and broader literary observation as the
text warrants.

This variability is a feature, not a deficiency. It mirrors the
way a knowledgeable reading companion would engage with a text ---
pausing where things are interesting, moving on where they are not.

\section*{Suitability for Other Languages and Learners}

The method is not specific to Polish. It is applicable to any
language for which a sufficiently capable LLM exists, and to any
learner who has a genuine motivation to read a specific text in
that language. The key conditions that make it work are:

\begin{itemize}
  \item a source text that the learner \textit{actually wants to read}
    --- ideally one unavailable in their native language, removing
    the option of simply reading a translation;
  \item willingness to proceed slowly and to treat the analysis as
    enjoyable in its own right, not merely as a means to an end;
  \item basic familiarity with \LaTeX{} if the learner wishes to
    maintain a compiled document, though the conversational
    analysis alone is valuable without any document production.
\end{itemize}

The method works particularly well for learners with existing
metalinguistic awareness --- those who have studied another
inflected language (Latin, Russian, German, classical Greek) and
already understand concepts like case, aspect, and agreement
abstractly. For such learners, the AI analysis illuminates the
\textit{Polish instantiation} of familiar concepts rather than
introducing entirely new ideas from scratch.

\section*{Transparency and Limitations}

The analyses in this document reflect the capabilities of the model
at the time of writing and should be treated as a knowledgeable
first-pass rather than authoritative scholarship. Occasional errors
in grammatical analysis or nuance are possible. Where the author
has Polish-speaking family or other expert resources available,
those remain the final arbiter of correctness.

The quoted Polish passages in this document are brief and used
solely for educational commentary. They do not constitute
reproduction of the source text and are consistent with fair use
principles for personal, non-commercial educational purposes.
The document is shared openly but the source book remains under
its original copyright.

\vspace{1.5em}
\noindent The full project, including all \LaTeX{} source files,
is available at:\\
\href{https://github.com/}{[GitHub repository URL to be added]}

\vspace{0.5em}
\noindent Readers are welcome to adapt the structure, commands,
and approach for their own language learning projects.
